\section{Abundance for fourfolds with nonzero Euler characteristic}
\label{sec:euler-route}

The metric theorem gives abundance for klt fourfolds when
the Euler characteristic is nonzero. A published nonvanishing
criterion uses this Euler hypothesis to produce the first
section; a second criterion then gives semiampleness. We state
both inputs in terms of the actual adjoint bundle on a resolution,
so that the same zero-Lelong metric supplies their analytic
hypotheses.

\begin{corollary}\label[corollary]{cor:euler-abundance}
Let $(X,\Delta)$ be a projective complex klt fourfold with effective
rational boundary and nef adjoint $D=K_X+\Delta$.
If $\chi(X,\cO_X)\ne0$, then $D$ is semiample.
\end{corollary}

This proof combines \Cref{thm:adjoint-metric} with published
nonvanishing and semiampleness criteria. It isolates the use of
$\chi(X,\cO_X)\ne0$: that condition enters the nonvanishing
criterion, whereas the metric theorem has no Euler-characteristic
hypothesis. Related work on multiplier-ideal approximations and
supercanonical currents studies this analytic setting
\cite[Theorems~A--B]{Lazic2024Metrics}.

\subsection*{The two algebraic inputs}

For a $\mathbb Q$-Cartier divisor $P$ on a normal projective variety,
the metric condition in \cite[Definition~2.12]{LazicPeternell2020}
requires a positive integer $r$, with $rP$ Cartier, and a resolution
$f:Z\to X$ carrying a singular metric on $f^*\mathcal O_X(rP)$ whose
curvature has the form
\begin{equation}\label{mm:eq:generalized-algebraic}
 T=T_0+\sum_j a_j[E_j],\qquad
 T_0\geq0,\quad \nu(T_0,z)=0\ (z\in Z),\quad a_j\in\mathbb Q_{\geq0}.
\end{equation}
The divisor sum is finite. This is called a metric with
\emph{generalized algebraic singularities}. The sum may be zero;
in particular, every semipositive metric with zero Lelong numbers
everywhere meets this definition.

\begin{theorem}[Lazi\'c--Peternell, fourfold nonvanishing]
\label{mm:thm:lp-fourfold}
Let $(X,\Delta)$ be a normal projective complex klt pair of dimension
four, with effective rational boundary, and put $D=K_X+\Delta$.
Suppose $D$ is pseudoeffective and $N$ is a nef $\mathbb Q$-Cartier
divisor. If both $D$ and $D+N$ satisfy the metric condition
\eqref{mm:eq:generalized-algebraic} and $\chi(X,\mathcal O_X)\ne0$,
then some rational $t_0>0$ satisfies
\[
 \kappa(X,D+tN)\geq0\qquad
 \text{for every rational }t\in[0,t_0].
\]
\end{theorem}

This is \cite[Corollary~D]{LazicPeternell2020}, the unconditional
four-dimensional consequence of their Theorem~C(i). Its interval includes
$t=0$. We will take $N=0$, so both metric assumptions concern $D$.

\begin{theorem}[Gongyo--Matsumura, fourfold semiampleness]
\label{mm:thm:gm-fourfold}
Let $(X,\Delta)$ be a normal projective complex klt pair of dimension
four, with effective rational boundary, and set $D=K_X+\Delta$.
Suppose that for a projective birational morphism $f:Z\to X$ with
$Z$ smooth and some positive integer $r$ with $rD$ Cartier, the bundle
$f^*\mathcal O_X(rD)$ carries a semipositive singular Hermitian metric
with zero Lelong numbers at every point of $Z$.
If $\kappa(X,D)\geq0$, then $D$ is semiample.
\end{theorem}

This is the form of \cite[Corollary~5.3]{GongyoMatsumura2017}
used here. Positivity of the Cartier multiple is also explicit in
their Theorem~5.1, from which that corollary follows. Neither of these
fourfold results requires $X$ to be $\mathbb Q$-factorial or
non-uniruled. Their dimension-three inputs are already incorporated
in the stated corollaries.

\subsection*{Application to the minimal metric}

\begin{proof}[Proof of Corollary~\ref{cor:euler-abundance}]
Choose a projective log resolution $\pi:Y\to X$, a positive integer
$r$ such that $rD$ is Cartier, and a metric $h_{\min}$ with minimal
singularities on $L=\pi^*D$. Nefness makes $D$ pseudoeffective,
so such a metric exists. By Theorem~\ref{thm:adjoint-metric},
\[
 \nu(h_{\min},y)=0\qquad(y\in Y).
\]
The tensor power $h_{\min}^{\otimes r}$ is a metric on the actual
line bundle $\pi^*\mathcal O_X(rD)$. Its local weights are $r\phi$,
so its curvature is semipositive and its Lelong numbers are
$r\nu(\phi,y)=0$.

Consequently, \eqref{mm:eq:generalized-algebraic} holds with
$T_0$ equal to this curvature current and with zero divisor sum.
Apply Theorem~\ref{mm:thm:lp-fourfold} with $N=0$ and $t=0$ to obtain
\[
 \kappa(X,D)\geq0.
\]
This is nonvanishing for $D$ itself: a positive Cartier multiple
of $D$ has a nonzero global section. It is therefore the Kodaira
dimension hypothesis of Theorem~\ref{mm:thm:gm-fourfold}. The same
resolution and the same metric $h_{\min}^{\otimes r}$ supply its
remaining hypothesis, and that theorem proves that $D$ is semiample.
\end{proof}

The argument imposes no restriction on the numerical dimension of $D$.
The Euler-characteristic hypothesis is used only to obtain the first
section; the interior-injectivity theorem of this article is not an
input to this particular application.
